\documentclass[11pt]{article}
\usepackage[a4paper,margin=1in]{geometry}
\usepackage{amsmath,amssymb,amsthm,mathtools,bm}
\usepackage{hyperref}
\usepackage{enumitem}
\newcommand{\R}{\mathbb{R}}
\newcommand{\X}{\mathcal{X}}
\newcommand{\A}{\mathcal{A}}
\newcommand{\Psafe}{\mathcal{P}_{\mathrm{safe}}}
\newcommand{\KL}{D_{\mathrm{KL}}}
\newcommand{\E}{\mathbb{E}}
\newtheorem{definition}{Definition}
\newtheorem{proposition}{Proposition}

\title{\textbf{A Mathematical Framework for Neural Privilege Separation}\\
\large Toward a Formal Theory of the Neural Firewall}
\author{NPS Research Project}
\date{\today}

\begin{document}
\maketitle

\begin{abstract}
This document develops a preliminary mathematical framework for Neural
Privilege Separation (NPS), with the goal of formalizing a neural firewall as
an internal security boundary within a transformer model. The central
hypothesis is that policy-relevant computation may admit a representation
that can be identified, causally manipulated, and ultimately protected from
attacker-controlled input while preserving legitimate model capabilities.
The framework deliberately separates mathematical definitions from empirical
claims: the existence of a useful policy representation and the feasibility
of the proposed guarantees remain open research questions.
\end{abstract}

\section{Research Goal}

The long-term objective of NPS is to treat policy enforcement as a property
of internal model computation rather than solely as an external guardrail.

The desired security property is analogous to privilege separation:
user-controlled computation should be able to access legitimate model
capabilities without being able to arbitrarily modify the internal state that
governs whether those capabilities may be used.

The mathematical objectives are to define:
\begin{enumerate}[label=\arabic*.]
\item a policy-relevant internal state,
\item a policy-safe region,
\item a mechanism for constraining internal computation,
\item an information-flow boundary,
\item a policy invariant, and
\item robustness against an explicit attacker model.
\end{enumerate}

\section{The Transformer as a Dynamical System}

Let a transformer $M_\theta$ have $L$ computational layers. For input $x$,
\begin{equation}
h_0=E(x),
\end{equation}
and
\begin{equation}
h_{\ell+1}=F_\ell(h_\ell),\qquad \ell=0,\ldots,L-1.
\end{equation}

The output distribution is
\begin{equation}
p_\theta(y\mid x)=\operatorname{softmax}(W_oh_L).
\end{equation}

Thus the model is a dynamical system
\begin{equation}
h_0\rightarrow h_1\rightarrow\cdots\rightarrow h_L
\rightarrow p_\theta(y\mid x).
\end{equation}

\section{Formalizing the Security Problem}

Conceptually decompose the input as
\begin{equation}
x=(u,t),
\end{equation}
where $u$ is user-controlled information and $t$ is task/environment
information.

A security violation is formulated as a trajectory problem:
\begin{equation}
x\rightarrow h_0\rightarrow h_1\rightarrow\cdots\rightarrow h_L
\end{equation}
that causes the model to enter a policy-violating computational state.

Let
\begin{equation}
\pi(h_\ell)\in[0,1]
\end{equation}
be a statistical estimate of policy-violation probability at state
$h_\ell$. This is initially a representation-reading quantity, not itself a
firewall.

\section{The Policy-Safe Set}

Define a policy-safe region
\begin{equation}
\Psafe\subseteq\R^d.
\end{equation}

A state is policy-safe when
\begin{equation}
h_\ell\in\Psafe.
\end{equation}

This does not assume that policy is represented by a single vector
$v_{\mathrm{policy}}\in\R^d$. The safe set could instead be a half-space,
subspace, manifold, or more general region.

\section{A Neural Firewall Operator}

Introduce
\begin{equation}
\mathcal{F}_\ell:\R^d\rightarrow\R^d.
\end{equation}

The protected transition is
\begin{equation}
h_{\ell+1}=F_\ell\!\left(\mathcal{F}_\ell(h_\ell)\right).
\end{equation}

The firewall should preserve already-safe states:
\begin{equation}
\mathcal{F}_\ell(h)=h,\qquad \forall h\in\Psafe,
\end{equation}
and map unsafe states into the safe region:
\begin{equation}
h\notin\Psafe\Longrightarrow
\mathcal{F}_\ell(h)\in\Psafe.
\end{equation}

\begin{definition}[Policy-preserving firewall]
A firewall operator $\mathcal F$ is policy-preserving if
\begin{equation}
\mathcal F(h)\in\Psafe,\qquad\forall h\in\R^d,
\end{equation}
and
\begin{equation}
\mathcal F(h)=h,\qquad\forall h\in\Psafe.
\end{equation}
\end{definition}

Safety alone is insufficient: the degenerate map $\mathcal F(h)=0$ could
destroy all useful behavior. Capability preservation is therefore required.

\section{Capability Preservation}

Let $D$ measure behavioral deviation between output distributions. A
capability-preserving firewall seeks
\begin{equation}
\mathcal F^*(h)
=
\arg\min_{z\in\Psafe}
D\!\left(p_\theta(\cdot\mid h),p_\theta(\cdot\mid z)\right).
\end{equation}

For example,
\begin{equation}
D(p,q)=\KL(p\Vert q).
\end{equation}

The intended interpretation is: change internal state as little as
necessary to restore policy compliance.

\section{The Neural Privilege Separation Hypothesis}

The central NPS hypothesis is that hidden computation may admit
\begin{equation}
\phi:\R^d\rightarrow
\mathcal U\times\mathcal C\times\mathcal P,
\end{equation}
with
\begin{equation}
\phi(h)=(u,c,p),
\end{equation}
where $u$ is user/task information, $c$ is capability information, and $p$
is policy information.

This decomposition is a hypothesis, not an established property of
transformers.

\section{The Privilege Boundary}

Suppose the internal transition can be approximated by
\begin{align}
u'&=f_u(u,c),\\
c'&=f_c(u,c,p),\\
p'&=f_p(c,p).
\end{align}

The critical property is
\begin{equation}
p'=f_p(c,p),
\end{equation}
rather than
\begin{equation}
p'=f_p(u,c,p).
\end{equation}

The intended information-flow restriction is therefore
\begin{equation}
U\not\rightarrow P
\end{equation}
within the protected policy transition, while legitimate flow such as
\begin{equation}
U\rightarrow C
\end{equation}
remains possible.

\section{Local Policy Isolation}

Define local policy sensitivity as
\begin{equation}
J=
\left\|
\frac{\partial p'}{\partial u}
\right\|.
\end{equation}

Approximate isolation requires
\begin{equation}
\boxed{
\left\|
\frac{\partial p'}{\partial u}
\right\|\leq\epsilon.
}
\end{equation}

An empirical finite-difference analogue is
\begin{equation}
\Delta_u=
\frac{
\left\|\phi_p(h(u+\delta))-\phi_p(h(u))\right\|
}{
\|\delta\|
}.
\end{equation}

\section{Global Policy Isolation}

Let $\A(x)$ be the set of allowed attacker transformations of $x$. Require
\begin{equation}
\boxed{
\sup_{x'\in\A(x)}
d_P\!\left(p(x'),p(x)\right)\leq\epsilon,
}
\end{equation}
where $d_P$ is a metric over policy representations.

This expresses the intended property that no permitted attacker transformation
can move protected policy state by more than a bounded amount.

\section{Policy Invariance}

Let $\Psafe$ be the safe policy set. Assume
\begin{equation}
p_0\in\Psafe
\end{equation}
and
\begin{equation}
p_\ell\in\Psafe
\Longrightarrow
p_{\ell+1}\in\Psafe.
\end{equation}

\begin{proposition}[Policy-state invariance]
Under the above conditions,
\begin{equation}
\boxed{
p_\ell\in\Psafe,\qquad\forall\ell.
}
\end{equation}
\end{proposition}

\begin{proof}
By induction. The base case is $p_0\in\Psafe$. If $p_\ell\in\Psafe$, the
assumed transition property gives $p_{\ell+1}\in\Psafe$. Therefore the
property holds for all layers.
\end{proof}

This is the mathematical analogue of protecting privileged state from being
driven outside its permitted region.

\section{Connecting Policy State to Behavior}

Let
\begin{equation}
V(y)\in[0,1]
\end{equation}
be a policy-violation function over outputs. Define
\begin{equation}
R(x)=
\E_{y\sim p_\theta(\cdot\mid x)}[V(y)].
\end{equation}

For a protected model $M_F$, the desired security condition is
\begin{equation}
\boxed{
\sup_{x\in\X_{\mathrm{adv}}}R_F(x)\leq\epsilon.
}
\end{equation}

\section{Capability Preservation Constraint}

Let $\X_{\mathrm{benign}}$ denote legitimate tasks. Require
\begin{equation}
\E_{x\sim\X_{\mathrm{benign}}}
D\!\left(p_F(\cdot\mid x),p_M(\cdot\mid x)\right)
\leq\delta.
\end{equation}

Thus the firewall must satisfy both security and capability preservation:
\begin{align}
\sup_{x\in\X_{\mathrm{adv}}}R_F(x)&\leq\epsilon,\\
\E_{x\sim\X_{\mathrm{benign}}}D(p_F,p_M)&\leq\delta.
\end{align}

\section{Recovery and Adversarial Robustness}

Let $\mathcal A$ denote an attacker class capable of optimizing prompts or
other permitted inputs. Define recovery risk
\begin{equation}
R_A(M_F)=
\max_{A\in\mathcal A}
V\!\left(A(M_F,x)\right).
\end{equation}

A robust firewall should satisfy
\begin{equation}
\boxed{
R_A(M_F)\leq\rho.
}
\end{equation}

The intended requirement is that the policy constraint remains effective under
adversarial optimization rather than only under ordinary prompting.

\section{Representation Identifiability}

The existence of a useful policy representation cannot simply be assumed.
Let $\Pi$ denote the underlying policy state and
$X_{\mathrm{surface}}$ denote surface-level prompt properties.

A desired information-theoretic condition is
\begin{equation}
I(P;\Pi)>0,
\end{equation}
while ideally
\begin{equation}
I(P;X_{\mathrm{surface}}\mid\Pi)\approx0.
\end{equation}

The first condition expresses policy information content; the second
expresses conditional independence from irrelevant surface features.
Causal meaning must still be established separately through intervention.

\section{Policy--User Independence}

A stronger privilege-separation objective is
\begin{equation}
\boxed{
I(P;U\mid\Pi)\approx0.
}
\end{equation}

The intended interpretation is that protected policy state should not contain
unnecessary user-controlled degrees of freedom that can be exploited to
rewrite the policy mechanism.

\section{Neural Firewall Optimization Problem}

Let $\theta_F$ denote firewall parameters and $\phi$ the representation map.
A preliminary optimization problem is
\begin{equation}
\boxed{
\begin{aligned}
\min_{\theta_F,\phi}\quad&
\mathcal L_{\mathrm{cap}}
+\lambda\mathcal L_{\mathrm{iso}}
\\
\text{s.t.}\quad&
\sup_{x\in\X_{\mathrm{adv}}}R_F(x)\leq\epsilon,
\\
&
\sup_{x'\in\A(x)}
d_P\!\left(p_F(x'),p_F(x)\right)\leq\eta,
\\
&
\sup_{A\in\mathcal A}R_A(M_F)\leq\rho.
\end{aligned}
}
\end{equation}

This is a research formulation, not a solved optimization problem.

\section{The Central Mathematical Question}

The unresolved question is whether a transformer admits a useful
\begin{equation}
\phi(h)=(u,c,p)
\end{equation}
for which the following can simultaneously hold:
\begin{align}
I(P;\Pi)&>0,\\
I(P;U\mid\Pi)&\approx0,\\
\left\|\frac{\partial p'}{\partial u}\right\|&\leq\epsilon,\\
\sup_{x'\in\A(x)}d_P(p(x'),p(x))&\leq\eta,\\
p_\ell\in\Psafe&\Rightarrow p_{\ell+1}\in\Psafe,\\
\sup_{x\in\X_{\mathrm{adv}}}R_F(x)&\leq\epsilon,\\
\E_{x\sim\X_{\mathrm{benign}}}D(p_F,p_M)&\leq\delta,\\
\sup_{A\in\mathcal A}R_A(M_F)&\leq\rho.
\end{align}

If such a representation and transition system can be constructed and
validated, NPS moves from activation steering toward an explicit theory of
privileged internal computation.

\section{Research Program}

\subsection{Stage 1: Representation}
Establish empirically whether a policy-relevant representation exists:
\begin{equation}
h\rightarrow p.
\end{equation}

\subsection{Stage 2: Causality}
Establish that intervention on $p$ changes policy-relevant behavior:
\begin{equation}
do(p\leftarrow p')\Longrightarrow Y\text{ changes}.
\end{equation}

\subsection{Stage 3: Isolation}
Measure whether attacker-controlled input can alter protected policy state:
\begin{equation}
\sup_{x'\in\A(x)}d_P(p(x'),p(x))\leq\epsilon.
\end{equation}

\subsection{Stage 4: Invariance}
Establish that the protected transition preserves the safe set:
\begin{equation}
p_\ell\in\Psafe\Longrightarrow p_{\ell+1}\in\Psafe.
\end{equation}

\subsection{Stage 5: Robustness}
Evaluate whether an explicit attacker can recover policy-violating behavior:
\begin{equation}
\sup_{A\in\mathcal A}V(A(M_F))\leq\epsilon.
\end{equation}

\section{Current Status}

The framework above consists of definitions, hypotheses, candidate
constraints, and desired theorems. The following remain unestablished:
\begin{itemize}
\item that a universal policy representation exists;
\item that user, capability, and policy states are cleanly separable;
\item that the policy state is low-dimensional;
\item that policy can be isolated from user-controlled computation;
\item that an invariant safe set exists for an arbitrary transformer;
\item that the proposed constraints can be jointly satisfied without
unacceptable capability loss.
\end{itemize}

\section{Conclusion}

The proposed mathematical view of a neural firewall is
\begin{equation}
\boxed{
\text{Protected policy state}
+
\text{constrained transition}
+
\text{capability-preserving projection}.
}
\end{equation}

The objective is not merely to identify a safety vector. It is to determine
whether policy-relevant computation can behave like privileged state: it
remains within a policy-safe region, preserves legitimate capabilities, and
resists modification or recovery through attacker-controlled inputs.

The strongest eventual result would therefore be a construction of a
protected policy state together with an explicit attacker model, measurable
isolation conditions, and a provable or bounded policy invariant.

\end{document}
